\section*{Chapitre \ref{ch:b1:gas-exchange} --- Les échanges gazeux respiratoires}

\begin{solution}{exo:b1:gas-exchange:1}
$\dot V = K(A/x)(P_1 - P_2)$~: une aire plus grande, une barrière plus
mince, une \omterm{def:b1:gas-exchange:partial}{pression partielle} externe plus élevée (ventilation) et une
pression interne plus basse (perfusion).
\end{solution}

\begin{solution}{exo:b1:gas-exchange:2}
Trente fois moins d'oxygène par litre (\qty{7}{} contre
\qty{210}{mL/L})~; huit cents fois plus dense et cinquante fois plus
visqueuse, donc coûteuse à pomper~; l'oxygène y diffuse dix mille fois
plus lentement, donc le milieu doit être amené jusqu'à la surface.
\end{solution}

\begin{solution}{exo:b1:gas-exchange:3}
Co-courant~: l'eau et le sang sortent tous deux à \qty{12}{kPa}~;
extraction $(21 - 12)/21 = \qty{43}{\%}$. \omterm{def:b1:gas-exchange:gill}{Contre-courant}~: l'eau sort
à \qty{4}{kPa}, le sang à \qty{18}{kPa}~; extraction $(21 - 4)/21 =
\qty{81}{\%}$.
\end{solution}

\begin{solution}{exo:b1:gas-exchange:4}
Dissous (\qty{7}{\%}), lié aux groupements aminés de l'hémoglobine
(\qty{23}{\%}), en bicarbonate dans le \omterm{def:b1:mammal-organization:compartments}{plasma} (\qty{70}{\%}).
\end{solution}

\begin{solution}{exo:b1:gas-exchange:5}
Minute \qty{6.75}{L/min}~; alvéolaire $15\times 0.3 = \qty{4.5}{L/min}$~;
prélèvement $4.5\times 0.07 = \qty{0.32}{L/min}$. Superficielle~: minute
\qty{6.75}{L/min} encore, mais alvéolaire $30\times 0.075 =
\qty{2.25}{L/min}$ et prélèvement \qty{0.16}{L/min} --- la moitié, pour
le même travail respiratoire~; l'espace mort se paie à chaque souffle.
\end{solution}

\begin{solution}{exo:b1:gas-exchange:6}
$2/(6\times 0.75) = \qty{0.44}{L}$ par heure, vingt-deux fois son
volume de \qty{20}{mL}.
\end{solution}

\begin{solution}{exo:b1:gas-exchange:7}
$A/x$ chute d'un facteur $2\times 4 = 8$~: \qty{31}{mL/min}. Pour
rétablir \qty{250}{}, la différence devrait être de \qty{64}{kPa} ---
impossible, puisque la pression alvéolaire ne peut dépasser
\qty{13}{kPa} dans l'air et que le sang veineux ne peut descendre
sous zéro~; le patient ne peut l'élever qu'en respirant de l'oxygène,
et pas d'un facteur huit même ainsi.
\end{solution}

\begin{solution}{exo:b1:gas-exchange:8}
L'\omterm{prop:b1:gas-exchange:transport}{anhydrase carbonique} est à l'intérieur de la \omterm{def:b1:cell-unit-of-life:cell}{cellule}~: le
bicarbonate y est fabriqué et s'y accumule~; il en sort par diffusion
le long de son gradient, via un échangeur qui admet un chlorure pour
chaque bicarbonate, ce qui garde les charges équilibrées. Sans cet
échange, le bicarbonate resterait à l'intérieur avec le partenaire de
son proton, ce qui élèverait l'osmolarité de la \omterm{def:b1:cell-unit-of-life:cell}{cellule}, et l'eau
entrerait~: la \omterm{def:b1:cell-unit-of-life:cell}{cellule} gonflerait dans les \omterm{def:b1:body-plans-tissues:tissue}{tissus} (elle gonfle un peu,
de fait).
\end{solution}

\begin{solution}{exo:b1:gas-exchange:9}
L'envie de respirer vient de la montée du $\mathrm{CO_2}$.
L'hyperventilation abaisse le $\mathrm{CO_2}$ très au-dessous de la
normale sans ajouter beaucoup d'oxygène (l'hémoglobine était déjà
saturée)~; pendant la plongée, l'oxygène tombe au niveau de la perte de
conscience avant que le $\mathrm{CO_2}$ ne soit remonté au niveau qui
force une inspiration. Sans hyperventilation, le $\mathrm{CO_2}$ atteint
ce niveau alors que l'oxygène est encore abondant, et le plongeur
remonte.
\end{solution}

\begin{solution}{exo:b1:gas-exchange:10}
$A = \pi(10^{-5})^2 = \qty{3.1e-10}{m^2}$~; $\Delta c = JL/DA =
10^{-9}\times 10^{-3}/(2\times 10^{-5}\times 3.1\times 10^{-10}) =
\qty{160}{mol/m^3}$ --- bien plus que les \qty{8.7}{} disponibles~: un
tube aussi étroit ne peut alimenter le muscle par diffusion sur
\qty{1}{mm}~; il faut un tube plus large (\qty{100}{\micro m}~:
$\Delta c = \qty{6.4}{mol/m^3}$, tout juste possible) ou une
ventilation active. À \qty{10}{mm}, l'exigence est décuplée encore.
Comme la longueur des tubes croît avec la taille du corps et la demande
avec le volume, la diffusion trachéenne plafonne la taille des insectes
à quelques centimètres, davantage seulement avec un pompage.
\end{solution}

\begin{solution}{exo:b1:gas-exchange:11}
Poisson~: eau dans un seul sens, sang opposé (\omterm{def:b1:gas-exchange:gill}{contre-courant}),
extraction \qty{80}{\%}. Oiseau~: air dans un seul sens à travers les
parabronches, sang à angle droit (courants croisés), extraction
intermédiaire mais supérieure à celle des mammifères, la \omterm{prop:b1:organism-environment:surfaces}{surface
d'échange} ne recevant jamais d'air vicié. Un \omterm{def:b1:gas-exchange:lung}{poumon} alternant mêle
l'air frais au gaz résiduel~: la \omterm{prop:b1:organism-environment:surfaces}{surface d'échange} ne voit au mieux que
le mélange alvéolaire, et le sang ne peut s'équilibrer qu'avec lui~:
quelle que soit la surface, le sang artériel ne peut dépasser la
pression alvéolaire, et la pression alvéolaire ne peut approcher celle
de l'air inspiré tant qu'un volume résiduel subsiste.
\end{solution}

\begin{solution}{exo:b1:gas-exchange:12}
Le $\mathrm{CO_2}$ est le meilleur signal parce que son niveau artériel
répond directement et fortement à la ventilation (divisez la
ventilation par deux et il double), parce qu'il est mesuré avec
précision par le pH des récepteurs centraux, et parce que l'oxygène,
assis sur le plateau de la courbe de l'hémoglobine, change peu sur le
domaine normal de ventilation~: un thermostat à
$\mathrm{CO_2}$ maintient l'oxygène correct par surcroît. Il échoue
quand l'oxygène chute indépendamment du $\mathrm{CO_2}$ --- en altitude,
où l'oxygène inspiré est bas mais la production de $\mathrm{CO_2}$ ne
l'est pas, de sorte que le signal $\mathrm{CO_2}$ dit «~respire moins~»
quand l'oxygène dit «~respire plus~»~; seuls les corpuscules
carotidiens, quand l'oxygène est très bas, le contredisent, et le
compromis (hyperventilation avec alcalose) demande des jours
d'ajustement rénal.
\end{solution}

\begin{solution}{pb:b1:gas-exchange:1}
\textbf{1.} Air $210/22.4 = \qty{9.4}{mmol/L}$~; eau $7/22.4 =
\qty{0.31}{mmol/L}$.
\textbf{2.} 30.
\textbf{3.} Air~: $1.2/9.4 = \qty{0.13}{g}$~; eau~: $1000/0.31 =
\qty{3200}{g}$ --- \num{25000} fois plus de masse.
\textbf{4.} $210/7 = \qty{30}{L}$ d'eau.
\textbf{5.} L'eau est trop lourde et trop pauvre en oxygène pour être
poussée en avant et en arrière~; elle doit être traversée une fois, ce
qui rend aussi le \omterm{def:b1:gas-exchange:gill}{contre-courant} possible~; l'air est assez bon marché
pour être déplacé deux fois.
\textbf{6.} Minute $12\times 0.5 = \qty{6}{L/min}$~; alvéolaire
$12\times 0.35 = \qty{4.2}{L/min}$.
\textbf{7.} $150/500 = \qty{30}{\%}$ de chaque souffle, et de la ventilation minute.
\textbf{8.} Délivré $4.2\times 0.21 = \qty{0.88}{L/min}$~; prélèvement
$4.2\times 0.07 = \qty{0.29}{L/min}$, proche de \qty{250}{mL/min}.
\textbf{9.} $0.25/(6\times 0.21) = \qty{20}{\%}$ (\qty{28}{\%} de ce
qui atteint les \omterm{def:b1:gas-exchange:lung}{alvéoles}).
\textbf{10.} $0.14\times(101 - 6.3) = \qty{13.3}{kPa}$.
\textbf{11.} $250/5 = \qty{50}{mL/L}$~; veineux \qty{150}{mL/L},
saturé à \qty{75}{\%}.
\textbf{12.} $3000/20 = \qty{150}{mL/L}$~; veineux \qty{50}{mL/L},
saturé à \qty{25}{\%}.
\textbf{13.} $50/60 = \qty{0.83}{mL/min}$.
\textbf{14.} $0.83/(7\times 0.8) = \qty{0.149}{L/min}$, \qty{8.9}{L}
par heure.
\textbf{15.} $0.83/(7\times 0.43) = \qty{0.28}{L/min}$~: le
\omterm{def:b1:gas-exchange:gill}{contre-courant} divise par deux l'eau à pomper, et le travail.
\textbf{16.} \qty{8.9}{kg} d'eau par heure~; l'humain déplace $6\times
60\times 1.2 = \qty{430}{g}$ d'air par heure --- un vingtième de la
masse, pour trois cents fois l'oxygène ($15\,000$ contre
\qty{50}{mL} par heure).
\textbf{17.} $100\times(0.95 - 0.15) = \qty{80}{mL/L}$~; débit cardiaque
$0.83/0.080 = \qty{10.4}{mL/min}$.
\textbf{18.} Eau/sang $149/10.4 = 14$~; air/sang $4.2/5 = 0.84$.
Le poisson doit faire passer quatorze volumes d'eau par volume de sang
parce que chacun contient si peu d'oxygène~; l'humain, environ un.
\textbf{19.} Prélèvement $\times 1.5$ à partir d'une eau qui en
contient $5.5/7 = 0.79$ fois autant~: ventilation
$\times 1.9$~; le coût de la respiration, à peu près proportionnel au
débit (ou davantage, puisqu'il croît plus vite que linéairement), passe
de \qty{10}{\%} à au moins \qty{13}{\%} d'un prélèvement plus grand ---
le poisson peine davantage à respirer en eau chaude.
\textbf{20.} Côté entrée d'eau~: $21 - 18 = \qty{3}{kPa}$~; côté sortie
d'eau~: $4 - 3 = \qty{1}{kPa}$~: positive aux deux bouts.
\textbf{21.} Côté entrée $21 - 3 = \qty{18}{kPa}$~; côté sortie $12 - 12 =
0$~: les deux fluides se sont équilibrés et il ne reste aucun gradient,
bien que l'eau contienne encore \qty{57}{\%} de son oxygène.
\textbf{22.} Au co-courant~: le sang s'équilibre avec un gaz alvéolaire
unique et bien mélangé et sort à sa pression~; il ne peut dépasser
\qty{13.3}{kPa} parce qu'il n'y a aucun gaz plus frais à rencontrer en
aval, comme il y en a dans le \omterm{def:b1:gas-exchange:gill}{contre-courant}.
\textbf{23.} Truite $2000/(2\times 10^{-4}) = \num{1e7}$ (cm$^2$ par cm)~;
humain $10^6/(5\times 10^{-5}) = \num{2e10}$~: rapport \num{2000}. Les
prélèvements $250/0.83 = 300$. La truite prend sept fois plus d'oxygène
par unité de $A/x$~: sa différence de pression moyenne à travers la
lamelle doit être plus grande --- ce que fournit le \omterm{def:b1:gas-exchange:gill}{contre-courant}, qui
garde un gradient sur toute la surface là où celui du \omterm{def:b1:gas-exchange:lung}{poumon} alternant
est faible vers la fin de l'équilibration.
\textbf{24.} Le $\mathrm{CO_2}$ est vingt-cinq fois plus soluble dans
l'eau que l'oxygène~: l'eau l'emporte aussi vite qu'il est produit, et
le $\mathrm{CO_2}$ du sang d'un poisson reste très bas.
\textbf{25.} \omterm{def:b1:gas-exchange:gill}{Branchie} environ \qty{80}{\%}, \omterm{def:b1:gas-exchange:lung}{poumon} environ \qty{20}{\%}
de l'oxygène du milieu inspiré~; le flux à sens unique de la \omterm{def:b1:gas-exchange:gill}{branchie}
et la disposition à \omterm{def:b1:gas-exchange:gill}{contre-courant} de l'eau et du sang font la
différence.
\end{solution}
